\title[Causality]{Causality: Lecture 11}
\author[Joris Mooij]{\\[5mm]\large Joris Mooij \& Patrick Forr\'e\\
\url{j.m.mooij@uva.nl}, \url{p.d.forre@uva.nl}}
\institute[UvA]{University of Amsterdam}
\date[2025-04-17]{\normalsize April 17th, 2025}
\maketitle

\part{$\sigma$-Separation and Acyclifications}
\frame{\partpage}

\begin{frame}{(Un)blockable non-colliders}
Remember: $\Sc^G(v) := \Anc^G(v) \cap \Desc^G(v)$.
\begin{Def}[\ref{def:unblockable_noncollider}]
  Let $G=(J,V,E,L)$ be a CDMG and $ \pi =\lp  v_0 \sus \cdots \sus v_n \rp $ a walk in $G$.
  We call a non-collider $v_k$ on $\pi$ an \emph{unblockable non-collider} on $\pi$ if it is not an end-node ($k \notin \{0,n\}$) and it only has outgoing edges on $\pi$ to nodes in the same strongly connected component of $G$. That is,
  it is one of the following patterns:
  \[\begin{array}{rccl}
    v_{k-1} \hut v_k \hus v_{k+1} & \text{ with } & v_{k-1} \in \Sc^G(v_{k})\\
    v_{k-1} \suh v_k \tuh v_{k+1} & \text{ with } & v_{k+1} \in \Sc^G(v_{k})\\
    v_{k-1} \hut v_k \tuh v_{k+1} & \text{ with } & v_{k-1} \in \Sc^G(v_k) \land v_{k+1} \in \Sc^G(v_k)
  \end{array}\]
  Otherwise, $v_k$ is called a \emph{blockable non-collider} on $\pi$.
  This means that $v_k$ is either an end-node ($k \in \lC 0,n\rC$) or it has at least one outgoing arrow 
  $v_k \tuh v_{k\pm1}$ pointing to a node $v_{k\pm1}$ that lies in a different strongly connected component than $v_k$, i.e.\ $v_{k\pm1} \notin \Sc^G(v_k)$.
\end{Def}
\end{frame}

\begin{frame}{Running example}
  \begin{Eg}
  \centering
  \begin{tikzpicture}
    \begin{scope}
      \node[ndout] (X1) at (0,0) {$X_1$};
      \node[ndout] (X2) at (1.5,0) {$X_2$};
      \node[ndout] (X3) at (3,0) {$X_3$};
      \node[ndout] (X4) at (0,-1.5) {$X_4$};
      \node[ndout] (X5) at (3,-1.5) {$X_5$};
      \node[ndout] (X6) at (1.5,-2.5) {$X_6$};
      \node[ndint] (X7) at (-1.5,-1.5) {$X_7$};
      \node[ndout] (X8) at (0,-3) {$X_8$};
      \draw[arout] (X1) -- (X4);
      \draw[arout] (X2) -- (X4);
      \draw[arout] (X2) -- (X3);
      \draw[arout] (X3) -- (X5);
      \draw[arout] (X4) -- (X5);
      \draw[arout] (X5) -- (X6);
      \draw[arout] (X6) -- (X4);
      \draw[arint] (X4) -- (X8);
      \draw[arout] (X7) -- (X8);
    \end{scope}
  \end{tikzpicture}
  \begin{itemize}
    \item Give an example of a walk with an unblockable noncollider;
    \item Give an example of a walk with an blockable noncollider.
  \end{itemize}
  \end{Eg}
\end{frame}

\begin{frame}{$\sigma$-blocked walks}
\begin{Def}[\ref{def:sigma_blocking}]
  Let $G=(J,V,E,L)$ be a CDMG and $ \pi =\lp  v_0 \sus \cdots \sus v_n \rp $ a walk in $G$,
  and $C \ins J \cup V$ a subset of nodes.
  We say that the walk $\pi$ is:
    \emph{$C$-$\sigma$-open} (or \emph{$\sigma$-open given $C$}) if and only if:
  \begin{enumerate}
      \item all colliders $v_k$ on $\pi$ are in $\Anc^G(C)$, \emph{and}:
    \item all blockable non-colliders $v_k$ on $\pi$ are \emph{not} in $C$.
      % all of its adjacent children, $v_{k\pm 1}$, on $\pi$ lie in its strongly connected component $\Sc^G(v_k)$. \Joris{the last part ``all of its neighboring children on $\pi$ lie in its strongly connected component'' is equivalent to ``unblockable'' (see Definition~\ref{def:unblockable_noncollider}); we might move that definition to here for convenience.}
  \end{enumerate}
  Otherwise, we say that the walk $\pi$ is
  \emph{$C$-$\sigma$-blocked}  (or \emph{$\sigma$-blocked given $C$}), that is:
  \begin{enumerate}
    \item there exists a collider $v_k$ on $\pi$ that is \emph{not} in $\Anc^G(C)$, \emph{or}:
    \item there exists a blockable non-collider $v_k$ on $\pi$ in $C$.
  \end{enumerate}  
  Note that unblockable non-colliders are always $C$-$\sigma$-open, regardless of the subset $C \ins V \cup J$.
\end{Def}
If $G$ is acyclic then all non-colliders are blockable, and then $\pi$ is $C$-$\sigma$-open iff $\pi$ is $C$-$d$-open.
\end{frame}

\begin{frame}{Running example}
  \begin{Eg}
  \centering
  \begin{tikzpicture}
    \begin{scope}
      \node[ndout] (X1) at (0,0) {$X_1$};
      \node[ndout] (X2) at (1.5,0) {$X_2$};
      \node[ndout] (X3) at (3,0) {$X_3$};
      \node[ndout] (X4) at (0,-1.5) {$X_4$};
      \node[ndout] (X5) at (3,-1.5) {$X_5$};
      \node[ndout] (X6) at (1.5,-2.5) {$X_6$};
      \node[ndint] (X7) at (-1.5,-1.5) {$X_7$};
      \node[ndout] (X8) at (0,-3) {$X_8$};
      \draw[arout] (X1) -- (X4);
      \draw[arout] (X2) -- (X4);
      \draw[arout] (X2) -- (X3);
      \draw[arout] (X3) -- (X5);
      \draw[arout] (X4) -- (X5);
      \draw[arout] (X5) -- (X6);
      \draw[arout] (X6) -- (X4);
      \draw[arint] (X4) -- (X8);
      \draw[arout] (X7) -- (X8);
    \end{scope}
  \end{tikzpicture}
  \begin{itemize}
    \item Give an example of a walk that is $\sigma$-blocked given $X_4$;
    \item Give an example of a walk that is $\sigma$-open given $X_4$;
  \end{itemize}
  \end{Eg}
\end{frame}

%\frame{\frametitle{$\sigma$-Blocking and $\sigma$-Separation}
%\begin{definition}
%  Let $G = \lt J,E,V,L \rt$ be a CDMG and $C \subseteq J \cup V$ a subset of nodes and $\pi$ a walk in $G$:
%  $$\pi = v_0 \sus \dots \sus v_n.$$
%  $\pi$ is called \emph{$C$-$\sigma$-blocked}, or \emph{$\sigma$-blocked by $C$}, iff
%  \begin{itemize}
%    \item $v_0 \in C$ or $v_n \in C$, or
%    \item it contains a collider $v_{k-1} \suh v_k \hus v_{k+1}$ with $v_k \not\in \Anc^G(C)$, or
%    \item it contains a non-collider
%      $v_{k-1} \sus v_k \tuh v_{k+1}$  with $v_k \in C$
%    \alert{and such that $v_{k+1} \notin \Sc^{G}(v_k)$}, or
%    \item it contains a non-collider
%      $v_{k-1} \hut v_k \sus v_{k+1}$ with $v_k \in C$
%    \alert{and such that $v_{k-1} \notin \Sc^{G}(v_k)$}.
%  \end{itemize}
%  We say that \emph{$A$ is $\sigma$-separated from $B$ by $C$}, denoted \emph{$A \Perp^{\alert{\sigma}}_G B \mid C$}, if every walk with one end node in $A$ and one end node in $B \cup J$ is $\sigma$-blocked by $C$.
%\end{definition}
%}

\begin{frame}{$\sigma$-Blocking and $\sigma$-Separation}
\begin{Def}[\ref{def:sigma_separation}]
    Let $G=(J,V,E,L)$ be a CDMG and $A,B,C \ins J \cup V$ (not necessarily disjoint) subset of nodes.
    We then say that:
    \begin{enumerate}
        \item \emph{$A$ is $\sigma$-separated from $B$ given $C$ in $G$}, in symbols:
    \[ A \Perp^\sigma_G B \given C, \]
      if every walk from a node in $A$ to a node in $B \alert{\cup J}$ is $C$-$\sigma$-blocked by $C$.
        \item If that property does not hold we will write:
            \[ A \nPerp^\sigma_G B \given C. \]
        \item We also define the special case:
                \[ A \Perp^\sigma_G B \qquad :\iff \qquad A \Perp^\sigma_GB \given \emptyset.\]
    \end{enumerate}
\end{Def}
\end{frame}

\begin{frame}{Running example}
  \begin{Eg}
  \centering
  \begin{tikzpicture}
    \begin{scope}
      \node[ndout] (X1) at (0,0) {$X_1$};
      \node[ndout] (X2) at (1.5,0) {$X_2$};
      \node[ndout] (X3) at (3,0) {$X_3$};
      \node[ndout] (X4) at (0,-1.5) {$X_4$};
      \node[ndout] (X5) at (3,-1.5) {$X_5$};
      \node[ndout] (X6) at (1.5,-2.5) {$X_6$};
      \node[ndint] (X7) at (-1.5,-1.5) {$X_7$};
      \node[ndout] (X8) at (0,-3) {$X_8$};
      \draw[arout] (X1) -- (X4);
      \draw[arout] (X2) -- (X4);
      \draw[arout] (X2) -- (X3);
      \draw[arout] (X3) -- (X5);
      \draw[arout] (X4) -- (X5);
      \draw[arout] (X5) -- (X6);
      \draw[arout] (X6) -- (X4);
      \draw[arint] (X4) -- (X8);
      \draw[arout] (X7) -- (X8);
    \end{scope}
  \end{tikzpicture}
  \begin{itemize}
    \item Give an example of a $\sigma$-separation given $X_4$;
    \item Give an example of a $\sigma$-connection given $X_4$.
  \end{itemize}
  \end{Eg}
\end{frame}

\frame{\frametitle{Example: $\sigma$-Separation}
\begin{example}
  \begin{tikzpicture}
      \node at (1,3.5) {Graph $G^+(M)$:};
      \node[ndout] (X1) at (0,0) {$X_1$};
      \node[ndout] (X2) at (1.5,0) {$X_2$};
      \node[ndout] (X3) at (1.5,1.5) {$X_3$};
      \node[ndout] (X4) at (0,1.5) {$X_4$};
      \node[ndlat] (W1) at (-1,-1) {$W_1$};
      \node[ndlat] (W2) at (2.5,-1) {$W_2$};
      \node[ndlat] (W3) at (2.5,2.5) {$W_3$};
      \node[ndlat] (W4) at (-1,2.5) {$W_4$};
      \draw[arout] (X1) edge (X2);
      \draw[arout] (X2) edge (X3);
      \draw[arout] (X3) edge (X4);
      \draw[arout] (X4) edge (X1);
      \draw[arout] (W1) edge (X1);
      \draw[arout] (W2) edge (X2);
      \draw[arout] (W3) edge (X3);
      \draw[arout] (W4) edge (X4);
      \node at (-3,3.5) {SCM $M$:};
      \node[text width=4cm] at (-3,1) {\small $X_1=X_4+W_1$\\
      $X_2=X_1 \cdot W_2$\\
      $X_3=X_2+W_3$\\
      $X_4=X_3\cdot W_4$\\
      \,\\
      $W_1,W_2,W_3,W_4$\\ independent};
      \node at (5,1.75) {is $X_1 \displaystyle\Perp^d_{G^+(M)} X_3 \given X_2, X_4$?};
%      \node at (5,1) {but};
      \node at (5,0.25) {is $X_1 \displaystyle\Perp^\sigma_{G^+(M)} X_3 \given X_2, X_4$?};
  \end{tikzpicture}
\end{example}

\bigskip
In general: No $\sigma$-separations (given any set) between nodes within the same strongly connected component.
}

\begin{frame}{Graphical Acyclifications: Definition}
\begin{Def}[\ref{def:cdg_acyclification}]
  Given a CDMG $G = ( J, V, E, L)$, we call a CADMG $G' = (J, V, E', L')$
  an \emph{acyclification of $G$} if
  \begin{enumerate}
    \item $G'$ is acyclic;
    \item $G'$ has the same input nodes $J$ and output nodes $V$ as $G$;
    \item for every pair of nodes $(i,j)$ such that $i \not\in\Sc^G(j)$:
      \begin{enumerate}
      \item $i \tuh j \in E'$ iff there exists a node $j' \in \Sc^G(j)$ such that $i \tuh j' \in E$;
      \item $i \huh j \in L'$ iff there exist nodes $i' \in \Sc^G(i), j' \in \Sc^G(j)$ such that $i' \huh j' \in L$;
      \end{enumerate}
  \item for every pair of distinct nodes $(i,j)$ such that $i \in\Sc^G(j)$: $i \tuh j \in E'$ or $i \hut j \in E'$ or $i \huh j \in L'$.
  \end{enumerate}
\end{Def}
\end{frame}

\begin{frame}{Graphical Acyclifications: Example}
  \centering
  \scalebox{0.8}{\begin{tikzpicture}
    \begin{scope}
      \node at (0.75,1) {CDMG $G$:};
      \node[ndout] (X1) at (0,0) {$X_1$};
      \node[ndout] (X2) at (1.5,0) {$X_2$};
      \node[ndout] (X3) at (3,0) {$X_3$};
      \node[ndout] (X4) at (0,-1.5) {$X_4$};
      \node[ndout] (X5) at (3,-1.5) {$X_5$};
      \node[ndout] (X6) at (1.5,-2.5) {$X_6$};
      \node[ndint] (X7) at (-1.5,-1.5) {$X_7$};
      \node[ndout] (X8) at (0,-3) {$X_8$};
      \draw[arout] (X1) -- (X4);
      \draw[arout] (X2) -- (X4);
      \draw[arout] (X2) -- (X3);
      \draw[arout] (X3) -- (X5);
      \draw[arout] (X4) -- (X5);
      \draw[arout] (X5) -- (X6);
      \draw[arout] (X6) -- (X4);
      \draw[arint] (X4) -- (X8);
      \draw[arout] (X7) -- (X8);
    \end{scope}
    \node at (0.75,-4) {Two acyclifications of $G$:};
    \begin{scope}[yshift=-5cm,xshift=-3.5cm]
      \node[ndout] (X1) at (0,0) {$X_1$};
      \node[ndout] (X2) at (1.5,0) {$X_2$};
      \node[ndout] (X3) at (3,0) {$X_3$};
      \node[ndout] (X4) at (0,-1.5) {$X_4$};
      \node[ndout] (X5) at (3,-1.5) {$X_5$};
      \node[ndout] (X6) at (1.5,-2.5) {$X_6$};
      \node[ndint] (X7) at (-1.5,-1.5) {$X_7$};
      \node[ndout] (X8) at (0,-3) {$X_8$};
    \uncover<2- |handout:0>{
      \draw[arout] (X1) -- (X4);
      \draw[arout] (X2) -- (X4);
      \draw[arout] (X2) -- (X3);
      \draw[arout] (X3) -- (X5);
      \draw[arout] (X7) -- (X8);
      \draw[arint] (X4) -- (X8);
      \draw[arout] (X4) -- (X5);
      \draw[arout] (X5) -- (X6);
      \draw[arout] (X4) -- (X6);
      \draw[arint] (X1) -- (X5);
      \draw[arint] (X1) -- (X6);
      \draw[arout] (X2) -- (X5);
      \draw[arout] (X2) -- (X6);
      \draw[arout] (X3) -- (X4);
      \draw[arout] (X3) -- (X6);
    }
    \end{scope}
    \begin{scope}[yshift=-5cm,xshift=3.5cm]
      \node[ndout] (X1) at (0,0) {$X_1$};
      \node[ndout] (X2) at (1.5,0) {$X_2$};
      \node[ndout] (X3) at (3,0) {$X_3$};
      \node[ndout] (X4) at (0,-1.5) {$X_4$};
      \node[ndout] (X5) at (3,-1.5) {$X_5$};
      \node[ndout] (X6) at (1.5,-2.5) {$X_6$};
      \node[ndint] (X7) at (-1.5,-1.5) {$X_7$};
      \node[ndout] (X8) at (0,-3) {$X_8$};
    \uncover<2- |handout:0>{
      \draw[arout] (X1) -- (X4);
      \draw[arout] (X2) -- (X4);
      \draw[arout] (X2) -- (X3);
      \draw[arout] (X3) -- (X5);
      \draw[arint] (X4) -- (X8);
      \draw[arout] (X7) -- (X8);
      \draw[arlat] (X4) -- (X5);
      \draw[arlat] (X5) -- (X6);
      \draw[arlat] (X4) -- (X6);
      \draw[arint] (X1) -- (X5);
      \draw[arint] (X1) -- (X6);
      \draw[arout] (X2) -- (X5);
      \draw[arout] (X2) -- (X6);
      \draw[arout] (X3) -- (X4);
      \draw[arout] (X3) -- (X6);
    }
    \end{scope}
  \end{tikzpicture}}
\end{frame}

\begin{comment}
\begin{frame}{Graphical Acyclifications: Example}
  \scalebox{0.8}{\begin{tikzpicture}
    \begin{scope}
      \node at (1.5,2) {Graph:};
      \node[varh] (W1) at (0,1.25) {$W_1$};
      \node[varh] (W2) at (3,1.25) {$W_2$};
      \node[varh] (W3) at (1,0.5) {$W_3$};
      \node[varh] (W4) at (2,0.5) {$W_4$};
      \node[varh] (W5) at (1,-2) {$W_5$};
      \node[varh] (W6) at (2,-2) {$W_6$};
      \node[var] (X1) at (0,0) {$X_1$};
      \node[var] (X2) at (3,0) {$X_2$};
      \node[var] (X3) at (0,-1) {$X_3$};
      \node[var] (X4) at (3,-1) {$X_4$};
      \node[var] (X5) at (0,-2) {$X_5$};
      \node[var] (X6) at (3,-2) {$X_6$};
      \draw[arrh] (W1) -- (X1);
      \draw[arrh] (W2) -- (X2);
      \draw[arrh] (W3) -- (X3);
      \draw[arrh] (W4) -- (X4);
      \draw[arrh] (W5) -- (X5);
      \draw[arrh] (W6) -- (X6);
      \draw[arr] (X1) -- (X3);
      \draw[arr] (X2) -- (X4);
      \draw[arr] (X3) -- (X5);
      \draw[arr] (X4) -- (X6);
      \draw[arr,bend left=20] (X3) edge (X4);
      \draw[arr,bend left=20] (X4) edge (X3);
    \end{scope}
    \begin{scope}[yshift=-5cm,xshift=-5cm]
      \node[varh] (W1) at (0,1.25) {$W_1$};
      \node[varh] (W2) at (3,1.25) {$W_2$};
      \node[varh] (W3) at (1,0.5) {$W_3$};
      \node[varh] (W4) at (2,0.5) {$W_4$};
      \node[varh] (W5) at (1,-2) {$W_5$};
      \node[varh] (W6) at (2,-2) {$W_6$};
      \node[var] (X1) at (0,0) {$X_1$};
      \node[var] (X2) at (3,0) {$X_2$};
      \node[var] (X3) at (0,-1) {$X_3$};
      \node[var] (X4) at (3,-1) {$X_4$};
      \node[var] (X5) at (0,-2) {$X_5$};
      \node[var] (X6) at (3,-2) {$X_6$};
      \draw[arrh] (W1) -- (X1);
      \draw[arrh] (W2) -- (X2);
      \draw[arrh] (W3) -- (X3);
      \draw[arrh] (W4) -- (X4);
      \draw[arrh] (W3) -- (X4);
      \draw[arrh] (W4) -- (X3);
      \draw[arrh] (W5) -- (X5);
      \draw[arrh] (W6) -- (X6);
      \draw[arr] (X1) -- (X3);
      \draw[arr] (X2) -- (X4);
      \draw[arr] (X1) edge (X4);
      \draw[arr] (X2) edge (X3);
      \draw[arr] (X3) -- (X4);
      \draw[arr] (X3) -- (X5);
      \draw[arr] (X4) -- (X6);
    \end{scope}
    \begin{scope}[yshift=-5cm]
      \node at (1.5,2) {Various acyclifications:};
      \node[varh] (W1) at (0,1.25) {$W_1$};
      \node[varh] (W2) at (3,1.25) {$W_2$};
      \node[varh] (W3) at (1,0.5) {$W_3$};
      \node[varh] (W4) at (2,0.5) {$W_4$};
      \node[varh] (W5) at (1,-2) {$W_5$};
      \node[varh] (W6) at (2,-2) {$W_6$};
      \node[var] (X1) at (0,0) {$X_1$};
      \node[var] (X2) at (3,0) {$X_2$};
      \node[var] (X3) at (0,-1) {$X_3$};
      \node[var] (X4) at (3,-1) {$X_4$};
      \node[var] (X5) at (0,-2) {$X_5$};
      \node[var] (X6) at (3,-2) {$X_6$};
      \draw[arrh] (W1) -- (X1);
      \draw[arrh] (W2) -- (X2);
      \draw[arrh] (W3) -- (X3);
      \draw[arrh] (W4) -- (X4);
      \draw[arrh] (W3) -- (X4);
      \draw[arrh] (W4) -- (X3);
      \draw[arrh] (W5) -- (X5);
      \draw[arrh] (W6) -- (X6);
      \draw[arr] (X1) -- (X3);
      \draw[arr] (X2) -- (X4);
      \draw[arr] (X1) edge (X4);
      \draw[arr] (X2) edge (X3);
      \draw[arr] (X4) -- (X3);
      \draw[arr] (X3) -- (X5);
      \draw[arr] (X4) -- (X6);
    \end{scope}
    \begin{scope}[yshift=-5cm,xshift=5cm]
      \node[varh] (W1) at (0,1.25) {$W_1$};
      \node[varh] (W2) at (3,1.25) {$W_2$};
      \node[varh] (W3) at (1,0.5) {$W_3$};
      \node[varh] (W4) at (2,0.5) {$W_4$};
      \node[varh] (W5) at (1,-2) {$W_5$};
      \node[varh] (W6) at (2,-2) {$W_6$};
      \node[var] (X1) at (0,0) {$X_1$};
      \node[var] (X2) at (3,0) {$X_2$};
      \node[var] (X3) at (0,-1) {$X_3$};
      \node[var] (X4) at (3,-1) {$X_4$};
      \node[var] (X5) at (0,-2) {$X_5$};
      \node[var] (X6) at (3,-2) {$X_6$};
      \draw[arrh] (W1) -- (X1);
      \draw[arrh] (W2) -- (X2);
      \draw[arrh] (W3) -- (X3);
      \draw[arrh] (W4) -- (X4);
      \draw[arrh] (W3) -- (X4);
      \draw[arrh] (W4) -- (X3);
      \draw[arrh] (W5) -- (X5);
      \draw[arrh] (W6) -- (X6);
      \draw[arr] (X1) -- (X3);
      \draw[arr] (X2) -- (X4);
      \draw[arr] (X1) edge (X4);
      \draw[arr] (X2) edge (X3);
      \draw[biarr] (X4) -- (X3);
      \draw[arr] (X3) -- (X5);
      \draw[arr] (X4) -- (X6);
    \end{scope}
  \end{tikzpicture}}
\end{frame}
\end{comment}

\begin{frame}{Graphical Acyclifications: Purpose}
The important property of acyclifications is that they can be used to express $\sigma$-separation in a (possibly cyclic) graph
in terms of $d$-separation in an acyclification.

\begin{Prp}[\ref{prp:sigma_sep_as_d_sep}]
  Let $G = \lt J, V, E, L \rt$ be a CDMG and $G' = \lt J, V, E', L' \rt$ an acyclification of $G$.
  Then for $A, B, C \subseteq V \cup J$ (not necessarily disjoint) subsets of nodes,
  $$A \Perp^\sigma_G B \given C \iff A \Perp^\sigma_{G'} B \given C \iff A \Perp^d_{G'} B \given C$$
\end{Prp}
\begin{proof}
  Requires constructing $\sigma$-open paths in $G'$ from those in $G$, and vice versa. 
  See lecture notes.
\end{proof}
\end{frame}

\begin{frame}{Acyclification of iSCM: Definition}
  We also define an acyclification operation on simple iSCMs:
  \begin{Def}[\ref{def:scm_acyclification}]
    Let $M = \lt \Sin, \Send, \Srnd, \CX, P, f \rt$ be a simple iSCM with causal graph $G(M)$.
    For each strongly connected component $C$ of $G(M)$, let 
    $g^{[C]} : \CX_{\Sin \cup (\Send \setminus C) \cup \Srnd} \to \CX_C$
    be a (essentially unique) solution function for $M$ w.r.t.\ $C$.
    Define $\tilde{f} : \CXJ \times \CXV \times \CXW \to \CXV$ by its components
    $$\tilde{f}_v(\xJ,\xV,\xW) = g^{[\Sc^{G(M)}(v)]}_v(\xJ,x_{\Send\setminus C},\xW)$$
    for $v \in V$.
    $M^\acy = \ltuple \Sin, \Send, \Srnd, \CX, P, \tilde{f} \rtuple$ is called an
    \emph{acyclification of $M$}.
  \end{Def}
  Acyclifications are unique up to equivalence.
\end{frame}

\begin{frame}{Acyclification of iSCM: Purpose}
  \begin{Prp}[\ref{prp:scm_acyclification_properties}]
    The acyclifications $M^\acy$ of a simple iSCM $M$ are acyclic and observably equivalent to $M$.
%    The acyclification $M^\acy$ of a simple iSCM $M$ is acyclic, observably equivalent to $M$, and $G(M^\acy) \subseteq G'$ for any acyclification $G'$ of $G(M)$.
  \end{Prp}
  \begin{proof}
    By example. For a formal proof, see lecture notes.
    \begin{wrapfigure}{r}{0.4\linewidth}
    \scalebox{0.8}{\begin{tikzpicture}
      \begin{scope}
        \node[varh] (W1) at (1,1) {$W_1$};
        \node[varh] (W2) at (2,1) {$W_2$};
        \node[varh] (W3) at (1,0) {$W_3$};
        \node[varh] (W4) at (2,0) {$W_4$};
        \node[var] (X1) at (0,0.2) {$X_1$};
        \node[var] (X2) at (3,0.2) {$X_2$};
        \node[var] (X3) at (0,-1) {$X_3$};
        \node[var] (X4) at (3,-1) {$X_4$};
        \draw[arrh] (W1) -- (X1);
        \draw[arrh] (W2) -- (X2);
        \draw[arrh] (W3) -- (X3);
        \draw[arrh] (W4) -- (X4);
        \draw[arr] (X1) -- (X3);
        \draw[arr] (X2) -- (X4);
        \draw[arr,bend left] (X3) edge (X4);
        \draw[arr,bend left] (X4) edge (X3);
      \end{scope}
      \begin{scope}[yshift=-3cm]
        \node[varh] (W1) at (1,1) {$W_1$};
        \node[varh] (W2) at (2,1) {$W_2$};
        \node[varh] (W3) at (1,0.2) {$W_3$};
        \node[varh] (W4) at (2,0.2) {$W_4$};
        \node[var,fill=red] (X1) at (0,0) {$X_1$};
        \node[var,fill=green] (X2) at (3,0) {$X_2$};
        \node[var,fill=blue] (X3) at (0,-1) {$X_3$};
        \node[var,fill=blue] (X4) at (3,-1) {$X_4$};
        \draw[arrh] (W1) -- (X1);
        \draw[arrh] (W2) -- (X2);
        \draw[arrh] (W3) -- (X3);
        \draw[arrh] (W4) -- (X4);
        \draw[arrh] (W3) -- (X4);
        \draw[arrh] (W4) -- (X3);
        \draw[arr] (X1) -- (X3);
        \draw[arr] (X2) -- (X4);
        \draw[arr] (X1) edge (X4);
        \draw[arr] (X2) edge (X3);
      \end{scope}
    \end{tikzpicture}}
  \end{wrapfigure}
  Structural equations of $M$:
  \begin{align*}
    X_1 &= W_1 \qquad X_3 = X_1 X_4 + W_3 \\
    X_2 &= W_2 \qquad X_4 = X_2 X_3 + W_4
  \end{align*}
  Structural equations of $M^\acy$:
  \begin{align*}
    &{\color{red}X_1} = W_1 \qquad {\color{blue}X_3} = \frac{W_3 + {\color{red}X_1} W_4}{1 - {\color{red}X_1} {\color{green}X_2}} \\
    &{\color{green}X_2} = W_2 \qquad {\color{blue}X_4} = \frac{{\color{green}X_2} W_3 + W_4}{1 - {\color{red}X_1} {\color{green}X_2}}
  \end{align*}
  \end{proof}
\end{frame}

\begin{frame}{Compatibility of Graphical and iSCM Acyclification}
  The iSCM notion of acyclification is compatible with the graphical notion of acyclification:
  \begin{Prp}[\ref{prp:acyclifications_compatible}]
    Let $M$ be a simple iSCM.
    Then:
    \begin{itemize}
      \item $G^+(M^\acy) \subseteq \tilde{G}^+$ for any graphical acyclification $\tilde{G}^+$ of $G^+(M)$, and
      \item $G(M^\acy) \subseteq \tilde{G}$ for any graphical acyclification $\tilde{G}$ of $G(M)$.
    \end{itemize}
  \end{Prp}
\end{frame}

\part{Global Markov Property for Simple iSCMs}
\frame{\partpage}

\begin{frame}{Global Markov Property for Simple iSCMs}
  With the help of the acyclification, we immediately derive a Markov property for simple iSCMs from the one for causal Bayesian networks:
  \begin{Cor}[\ref{cor:gmp_simple_scm}]
    Let $M = \lt \Sin, \Send, \Srnd, \CX, P, f \rt$ be a simple iSCM with graph $G^+(M)$, causal graph $G(M)$
    and Markov kernel $P_{M}(X_V,X_W \given \doit(X_J))$.
    Then for all $A, B, C \ins J \cup V \cup W$ (not necessarily disjoint) we have the implication:
    $$A \Perp^\sigma_{G^+(M)} B \given C \implies X_A \Indep_{P_{M}(X_V,X_W \given \doit(\XJ))} X_B \given X_C$$
    and for all $A, B, C \ins J \cup V$:
    $$A \Perp^\sigma_{G(M)} B \given C \implies X_A \Indep_{P_{M}(X_V \given \doit(\XJ))} X_B \given X_C.$$
%    If one wants to make the implicit dependence on $J$ more explicit one can equivalently also write:
%    $$A \Perp^\sigma_{G^+(M)} J \cup B \given C \implies X_A \Indep_{P_{M}(X_V,X_W \given \doit(\XJ))} X_J,X_B \given X_C.$$
  \end{Cor}
\end{frame}

\begin{frame}{Proof of Global Markov Property for Simple iSCMs}
\begin{proof}
  \setlength\abovedisplayskip{10pt}
  \setlength\belowdisplayskip{0pt}
  Choose an acyclification $\tilde{G}^+$ of $G^+(M)$. Then:
  \begin{align*}
    A \Perp^\sigma_{G^+(M)} B \given C &\stackrel{1}{\iff} A \Perp^d_{\tilde{G}^+} B \given C \\
    &\stackrel{2}{\implies} A \Perp^d_{G^+(M^\acy)} B \given C \\
    &\stackrel{3}{\implies} X_A \Indep_{P_{M^{\acy}}(X_V,X_W \given \doit(\XJ))} X_B \given X_C \\
    &\stackrel{4}{\iff} X_A \Indep_{P_{M}(X_V,X_W \given \doit(\XJ))} X_B \given X_C.
  \end{align*}
  \begin{enumerate}
    \item $\tilde{G}^+$ is an acyclification of $G^+(M)$; 
    \item $G^+(M^\acy) \subseteq \tilde{G}^+$, removing edges preserves $d$-separations; 
    \item the GMP for $M^\acy$ interpreted as a causal Bayesian network;
    \item $M^\acy$ is observably equivalent to $M$.
  \end{enumerate}
\end{proof}
\end{frame}

\part{Do-Calculus for Simple iSCMs}
\frame{\partpage}

\begin{frame}{Hard Interventions Via Intervention Variables}
  \begin{Def}[\ref{def:scm_intervention_variables}]
  Let $M = \lt \Sin, \Send, \Srnd, \CX, P, f \rt$ be an iSCM, and $B \subseteq V \cup J$.
  For each $b \in B \cap V$, we introduce an additional `intervention variable' $I_b$ as exogenous input variable; for $j \in B \cap J$, we just set $I_j := j$.
  We denote $I_B := (I_b)_{b \in B}$.
  We define an extended iSCM 
  \vskip-2mm
  $$\textstyle M_{\doit(I_B)} := \lt \Sin \cup I_B, \Send, \Srnd, \CX \times \prod_{b \in B \cap V} \CX_{I_b}, P, \tilde{f} \rt$$
  with $\CX_{I_b} := \CX_b \,\dot\cup\, \{\star\}$, with causal mechanism $\tilde{f}$ with components
  \vskip-6mm
  \begin{align*}
    v \in V \cap B: \qquad &\tilde{f}_v(x_{V \cup J \cup W \cup I_B}) := \begin{cases}
      f_v(x_{V \cup J \cup W}) & x_{I_v} = \star \\
      x_{I_v} & x_{I_v} \in \CX_v
    \end{cases}\\
    v \in V \sm B: \qquad & \tilde{f}_v(x_{V \cup J \cup W \cup I_B}) := f_v(x_{V \cup J \cup W})
  \end{align*}
  \vskip-2mm
  For $b \in B \cap V$, $x_{I_b} = \star$ encodes that there is no intervention on $b$, while $x_{I_b} \ne \star$ encodes that the hard intervention $\doit(X_b = x_{I_b})$ is performed.
\end{Def}
Purpose: jointly encode $\{M_{\doit(X_C)} : C \subseteq B\}$ into a single iSCM $M_{\doit(I_B)}$.
\end{frame}

%\begin{frame}{Notation}
%  \begin{enumerate}
%    \item Let $M = \lt \Sin, \Send, \Srnd, \CX, P, f \rt$ be a simple iSCM with causal graph $G = G(M)$;
%    \item To model hard interventions $\doit(B)$ for $B \subseteq V \cup W$, introduce \emph{intervention variable} $I_B$;
%    \item Extended iSCM $\tilde{M} = \lt \Sin \cup \{I_B\}, \Send, \Srnd, \CX \times (\CX_B \dot\cup \{\star\}), P, \tilde{f} \rt$ with causal mechanism
%$$\tilde{f}(x,x_{I_B}) := \begin{cases}
%  f(x) & x_{I_B} = \star \\
%  (f_{\setminus B}(x), x_{I_B}) & x_{I_B} \in \CX_B
%\end{cases}$$
%    \item $x_{I_B} = \star$ encodes that there is no intervention on $B$,\\
%$x_{I_B} \ne \star$ encodes that the hard intervention $\doit(X_B = x_{I_B})$ is performed.
%\item The graph of $\tilde{M}$, $G_{\doit(I_B)}$, has an additional variable $I_B$ with no incoming edges
%  and outgoing edges to all $b \in B$.
%  \end{enumerate}
%%In the do-calculus, we make use of the extended graph $G_{\doit(I_B)}$ to test the separation statement, while the conclusion
%%about the existence and identity of particular Markov kernels concerns those of the original iSCM $M$.
%\end{frame}

\begin{frame}{Reminder: Notation of Markov kernels}
  For a simple iSCM $M$:
  \begin{enumerate}
    \item Its \emph{Markov kernel} is denoted $P_{M}(X_V \mid \doit(X_J))$.
    \item For a subset $T \subseteq V$, we write the \emph{intervened Markov kernel}
      $$P_{M}(X_{V\setminus T} \mid \doit(X_J,X_T)) = P_{M_{\doit(T)}}(X_{V\setminus T} \mid \doit(X_J,X_T)).$$
    \item By conditioning on a subset $S \subseteq V \sm T$, we obtain the conditional intervened Markov kernel
      $$P_{M}(X_{V\sm (T \cup S)} \mid \doit(X_J,X_T), X_S).$$
    \item Assume that we have $\sigma$-finite reference measures $\mu_v$ on $\Xcal_v$ for every $v \in V$
      and put $\mu_F := \bigotimes_{v \in F} \mu_v$ for $F \ins V$.
  \end{enumerate}
\end{frame}

%\begin{frame}{Do-Calculus for simple iSCMs (on one slide!)}
%  Let $A,B,C \ins V \cup W$ and $D \ins V \cup W \cup J$ be such that $A,B,C,D$ are pairwise disjoint.
%  \begin{enumerate}
%    \item \emph{Insertion/deletion of observation}: 
%      $A \Perp^\sigma_{G_{\doit(D)}} B \given C \cup D$ implies
%      \[ P_{M}\lp X_A|X_B,X_C,\doit(X_{D\cup J})\rp = P_{M}\lp X_A|X_C,\doit(X_D)\rp.\]
%    \item \emph{Action/observation exchange}: $A \Perp^\sigma_{G_{\doit(I_B,D)}} I_B \given B \cup C \cup D$ implies
%      \[ P_{M}\lp X_A|\doit(X_B),X_C,\doit(X_D)\rp  =  P_{M}\lp X_A|X_B,X_C,\doit(X_D)\rp.\]
%    \item \emph{Insertion/deletion of action}: $A \Perp^\sigma_{G_{\doit(I_B,D)}} I_B \given C  \cup D$ implies
%      \[ P_{M}\lp X_A|\doit(X_B),X_C,\doit(X_{D\cup J})\rp = P_{M}\lp X_A|X_C,\doit(X_{D})\rp.\]
%  \end{enumerate}
%  Note: all kernels should be interpreted as $\CX_{J \cup \dots} \dshto \dots$, but some are constant in $x_J$!
%\end{frame}

\begin{frame}{Do-Calculus for simple iSCMs, part I}
  \begin{Thm}[\ref{thm:as-do-calculus-scms-simplified}]
    Let $M$ be a simple iSCM with graph $G = G(M)$.
    Let $A,B,C \ins V$ and $D \ins V \cup J$ be such that $A,B,C,D$ are pairwise disjoint.
  \small
  \begin{enumerate}
      \item \emph{Insertion/deletion of observation}: if $J \ins D$, 
       \begin{align*} 
           A &\Perp^\sigma_{G_{\doit(D)}} B \given C \cup D, &
           \mu_{B \cup C} &\ll P_M(X_B,X_C|\doit(X_{D}))\ll \mu_{B \cup C},
       \end{align*}
          \[ \implies P_M(X_A|X_B,X_C,\doit(X_{D})) = P_M(X_A|X_C,\doit(X_{D})) \quad \mu_{B \cup C}\text{-a.s.}.\]
     \item \emph{Action/observation exchange}: if $J \ins D$,
       \begin{align*} 
           A & \Perp^\sigma_{G_{\doit(I_B,D)}} I_B \given B \cup C \cup D, &          
                  \mu_{B \cup C} &\ll  P_M(X_B,X_C|\doit(X_{D})) \ll \mu_{B \cup C},\\
               &&   \mu_C &\ll  P_M(X_C|\doit(X_B,X_{D})) \ll \mu_C,
       \end{align*}
       \[ \implies P_M(X_A|X_B,X_C,\doit(X_{D}))  = P_M(X_A|\doit(X_B),X_C,\doit(X_{D})) \quad \mu_{B \cup C}\text{-a.s.}.\]
    \end{enumerate}
\end{Thm}
\end{frame}

\begin{frame}{Do-Calculus for simple iSCMs, part II}
  \begin{Thm}[\ref{thm:as-do-calculus-scms-simplified}, continued]
  \small
  \begin{enumerate}
     \setcounter{enumi}{2}
     \item \emph{Insertion/deletion of action}: if $J \ins D$,
       \begin{align*} 
           A &\Perp^\sigma_{G_{\doit(I_B,D)}} I_B \given C  \cup D, &          
           \mu_C &\ll P_M(X_C|\doit(X_B,X_{D})) \ll \mu_C, \\
                 && \mu_C &\ll P_M(X_C|\doit(X_{D})) \ll \mu_C,
       \end{align*}
       \[ \implies  P_M(X_A|\doit(X_B),X_C,\doit(X_{D})) = P_M(X_A|X_C,\doit(X_{D})) \quad \mu_C\text{-a.s.}.\]
     \item \emph{Deletion of input}: 
         If 
         \begin{align*} 
           A &\Perp^\sigma_{G_{\doit(D)}} J \given C \cup D, &
             \mu_C &\ll P_M(X_C|\doit(X_{D\cup J})) \ll \mu_C,
         \end{align*}
         then there exists a Markov kernel $P_M(X_A|X_C,\doit(X_{D},\cancel{X_{J\sm D}}))$ such that:
         \[ P_M(X_A|X_C,\doit(X_{D},\cancel{X_{J\sm D}})) = P_M(X_A|X_C,\doit(X_{D \cup J})) \quad \mu_C\text{-a.s.}. \]
    \end{enumerate}
\end{Thm}
\end{frame}

%\begin{frame}{Proof of Do-Calculus for Simple iSCMs}
%  \begin{proof}
%  The proof is analogous to that of Corollary~\ref{cor:as-do-calculus-2}, except that it applies the global Markov property for simple iSCMs, Corollary~\ref{cor:gmp_simple_scm}, instead of the one for causal Bayesian networks, Theorem~\ref{thm-gmp-mCBN}.
%  \end{proof}
%\end{frame}

\begin{frame}{Adjustment formula for simple iSCMs}
Same notation as for do-calculus.
  \begin{block}{Problem setting (special case of ``Causal domain adaptation'')}
We want to estimate the \emph{conditional causal effect}
$P_{M}(X_A|X_C,\doit(X_B,X_D))$
from
$P_{M}(X_A,X_B,X_F|X_C,\doit(X_D)).$
\end{block}
The following (pairwise disjoint) index sets have the following roles:\\
\begin{description}\small
  \item[$A\ins V$]: the outcome variables of interest.
  \item[$B\ins V$]: the treatment or intervention variables.
  \item[$C\ins V$]: general conditional (context) variables under which the data was collected.
  \item[$J\ins D\ins V \cup J$]: general interventional (context) variables that were set by the experimenter, with $J \ins D$.
  \item[$F_0\ins V$]: core adjustment variables, i.e.\ features that were measured.
  \item[$F_1\ins V$]: additional measured adjustment variables, with $F = F_0 \cup F_1$.
%  \item[$H$]: additional unobserved variables.
\end{description}
\end{frame}

\begin{frame}{General adjustment formula for simple iSCMs}
  \begin{Thm}[\ref{thm:gen-adjustment-scms}]
    \small
Assume that all the following $\sigma$-separations hold in the graph $G_{\doit(I_B,D)}$:
    \begin{align*}
        F_0 & \Perp^\sigma_{G_{\doit(I_B,D)}} I_B \given (C \cup D), \\
        A & \Perp^\sigma_{G_{\doit(I_B,D)}} (F_1 \cup I_B) \given (B \cup F_0 \cup C \cup D).
%                      H &\Perp^\sigma_{G_{\doit(I_B,D)}} B \given (F \cup C \cup  I_B \cup D).
    \end{align*}
Further assume that we have reference measures $\mu_v$ on $\Xcal_v$, $v \in V$, such that:
  \begin{align*}
      \mu_{B \cup C \cup F} &\ll  P(X_B,X_C,X_F|\doit(X_D)) \ll \mu_{B \cup C \cup F},\\
      \mu_{C \cup F} &\ll  P(X_C,X_F|\doit(X_B,X_D)) \ll \mu_{C \cup F}.
  \end{align*}
Then we have the $\mu_{B \cup C}\text{-a.s.}$ adjustment formula:
      \[ P_{M}(X_A|X_C,\doit(X_B,X_D)) = P_{M}(X_A|X_B,X_C,X_F,\doit(X_D)) \circ P_{M}(X_F|X_C,\doit(X_D)) \]
\end{Thm}
\end{frame}

\begin{frame}{Proofs}
  \begin{proof}
    All proofs proceed the same as for L-CBNs, except that all $d$-separations need to be replaced with $\sigma$-separations, and we invoke the GMP for simple iSCMs rather than the GMP for L-CBNs.
  \end{proof}

  \bigskip
  Similarly:
  \begin{itemize}
    \item The more detailed statements of the do-calculus (Proposition~\ref{prp:do-calculus} and Theorem~\ref{thm:as-do-calculus-1}) are also valid for simple iSCMs;
    \item The interventional backdoor covariate adjustment formula (Corollary~\ref{cor:backdoor-int}) is also valid for simple iSCMs;
  \end{itemize}
  where we just need to replace $\Perp^d$ by $\Perp^\sigma$.
  
  \bigskip
  \begin{block}{Caveat}
  The literature often fails to mention the strict positivity assumptions, even though without sufficient positivity, the various backdoor criteria may not hold
    (e.g.~Example~\ref{eg:counterexample_nonpositive_backdoor_adjustment})!
  \end{block}
\end{frame}


\part{Bounds on causal effects}
\frame{\partpage}

\begin{frame}{Bounding non-identifiable causal effects}

  If causal effects cannot be identified from the observable distribution, we may still be able to derive informative \textbf{bounds}.

  \bigskip
    \begin{tikzpicture}
      \begin{scope}[xshift=0]
        \node[ndout] (X1) at (0,0) {$X_1$};
        \node[ndout] (X2) at (3,0) {$X_2$};
        \node[ndout] (X3) at (1.5,2.5) {$X_3$};
%        \node[ndlat] (E1) at (0,1.5) {$X_A$};
        \draw[arout] (X1) -- (X2);
        \draw[arout] (X3) -- (X1);
        \draw[arout] (X3) -- (X2);
        \node at (-1,2.5) {$G(M)$:};
      \end{scope}
      \begin{scope}[xshift=6cm]
        \node[ndout] (X1) at (0,0) {$X_1$};
        \node[ndout] (X2) at (3,0) {$X_2$};
        \node[ndout] (X3) at (1.5,2.5) {$X_3$};
%        \node[ndlat] (E1) at (0,1.5) {$X_A$};
        \draw[arout] (X1) -- (X2);
        \draw[arout, bend left] (X2) edge (X1);
        \draw[arout] (X3) -- (X1);
        \draw[arout] (X3) -- (X2);
        \draw[arout, bend left] (X1) edge (X3);
        \draw[arout, bend right] (X2) edge (X3);
        \draw[huh, bend right] (X2) edge (X1);
        \draw[huh, bend left] (X3) edge (X1);
        \draw[huh, bend left] (X2) edge (X3);
        \node at (-1,2.5) {$G(\tilde{M})$:};
      \end{scope}
    \end{tikzpicture}

  \bigskip
    \begin{itemize}
      \item $P_M(X_2 \given \doit(X_1), X_3)$ is identifiable (under positivity assumptions) from $P_M(X_1,X_2,X_3)$.% via adjustment on $X_3$. 
      \item $P_{\tilde{M}}(X_2 \given \doit(X_1), X_3)$ is not identifiable from $P_{\tilde{M}}(X_1,X_2,X_3)$, but we can still bound it using the natural bounds (if both $X_1$ and $X_3$ are discrete).
    \end{itemize}
\end{frame}

\begin{frame}{Natural bounds}
\cite{ManskiNagin1998} proved the following `natural' bounds. We point out here that they also hold for simple iSCMs.
\begin{Thm}[\ref{thm:natural_bounds}, Natural bounds on causal effect]
  Let $M$ be a simple iSCM with endogenous variables $V \supseteq \{1,2,3\}$ and no exogenous input variables.
  Assume that $\CX_1$ is discrete.
  Then
  \begin{equation*}\begin{split}
    &P_M(X_1=a,X_2 \in B|X_3 \in C) \\
    &\quad\le P_M(X_2\in B|X_3 \in C, \doit(X_1=a)) \\
    &\quad\le P_M(X_1=a,X_2\in B|X_3 \in C) + P_M(X_1\ne a|X_3 \in C).
  \end{split}\end{equation*}
  for any $a\in\C{X}_1$, measurable $B \subseteq \C{X}_2$, and measurable $C \ins \C{X}_3$ with $P_M(X_3 \in C) > 0$.
\end{Thm}
\end{frame}

\begin{frame}{Consistency of solutions of simple iSCMs}
\begin{Prp}[\ref{prp:simple_scm_consistency}]
  Let $M = \lt \Sin, \Send, \Srnd, \CX, P, f \rt$ be a simple iSCM.
  Let $g : \CXJ \times \CXW \to \CXV$ be a solution function of $M$.
  Let $V = A \dcup B$ be a partition of the endogenous variables of $M$,
  and let $g^{[B]} : \CX_A \times \CXJ \times \CXW \to \CX_B$ be a partial solution function of $M$ w.r.t.\ $B$.
  Then
  $$g_B(x_J,x_W) = g^{[B]}(g_A(x_J,x_W),x_J,x_W)\quad M\as.$$
\end{Prp}

\bigskip
\begin{block}{Remark}
  This proposition shows that the ``consistency assumption'' commonly made in the potential outcomes framework
  holds true for potential outcomes defined by simple iSCMs.
\end{block}
\end{frame}

\begin{frame}{Natural bounds: proof}
\begin{proof}
  We use consistency $X_2=X_2^{\doit(X_1=X_1)}$ and elementary probability theory:
    \begin{align*}
      & P_M(X_1=a,X_2\in B|X_3 \in C) \\
      & = P_M(X_1=a,X_2^{\doit(X_1=a)}\in B|X_3 \in C) \\
      & \le P_M(X_2^{\doit(X_1=a)}\in B|X_3 \in C) \\
      & = P_M(X_1=a,X_2^{\doit(X_1=a)}\in B|X_3 \in C) \\
      & + P_M(X_1\ne a,X_2^{\doit(X_2=a)}\in B|X_3 \in C) \\
      & \le P_M(X_1=a,X_2\in B|X_3 \in C) + P_M(X_1 \ne a|X_3 \in C)
    \end{align*}
  The statement follows since
  $$ P_M(X_2\in B|X_3 \in C, \doit(X_1=a)) = P_M(X_2^{\doit(X_1=a)}\in B|X_3 \in C).$$
\end{proof}
\end{frame}

\begin{frame}{Natural bounds: real-valued outcome $X_2$}
If one has a priori knowledge about the range of $X_2$ (in case it is real-valued), one can also derive a bound on the expected result of an intervention \cite{ManskiPepper2013}.
\begin{Cor}[\ref{cor:natural_bounds}]
  In the situation of Theorem~\ref{thm:natural_bounds}, suppose that $\CX_2 = [\alpha,\beta]$ and $0 < P_M(X_1=a|X_3\in C) < 1$. Then
  \begin{equation*}\begin{split}
    & P_M(X_1=a|X_3\in C) \Exp_M(X_2|X_1=a,X_3\in C) + \alpha P_M(X_1\ne a|X_3\in C) \\
    {}\le {}& \Exp_M(X_2|X_3\in C,\doit(X_1=a)) \\
    {}\le {}& P_M(X_1=a|X_3\in C) \Exp_M(X_2|X_1=a,X_3\in C) + \beta P_M(X_1\ne a|X_3\in C).
  \end{split}\end{equation*}
\end{Cor}
\end{frame}

\begin{frame}{Discussion}
  \begin{itemize}
    \item The natural bound can be shown to be tight. 
    \item No assumptions regarding the causal relations between the three endogenous variables are necessary.
    \item We can bound the causal effect of $X_1$ on $X_2$ in the presence of both confounding and cycles.
    \item There exists no analogous bound in case $X_1$ is real-valued.
  \end{itemize}
\end{frame}
